\subsection{Admissibility Envelope by Sector: $\mu_\rho$, $\lambda_\rho$, $A_\rho^{\max}$}
  \label{subsec:admissibility-envelope-by-sector}

  For a finite group $G$ with symmetric generating set $S$ of size $|S|$, the
  Laplacian eigenvalue in irreducible sector $\rho$ is
  \begin{equation}
    \lambda_\rho \;=\; |S| - \frac{1}{\dim\rho}\sum_{s\in S}\chi_\rho(s)
    \;\equiv\; |S|\,\mu_\rho,
    \qquad
    \mu_\rho \;=\; 1 - \frac{1}{|S|\dim\rho}\sum_{s\in S}\chi_\rho(s).
  \end{equation}
  The admissibility envelope (Eq.~48 of the main text) then gives
  \begin{equation}
    A_\rho^{\max} \;=\; \frac{c_\chi}{\sqrt{|S|\,\mu_\rho}}.
  \end{equation}

  For $Q_8$ with $S=\{\pm i,\pm j,\pm k\}$, $|S|=6$:

  \begin{center}
    \begin{tabular}{lcccc}
      \hline
      Sector $\rho$ & $\dim\rho$ & $\mu_\rho$ & $\lambda_\rho$ & $A_\rho^{\max}/c_\chi$ \\
      \hline
      $\rho_1$ (trivial) & 1 & 0 & 0 & $\infty$ (zero mode) \\
      $\rho_2,\rho_3,\rho_4$ (abelian) & 1 & $4/3$ & 8 & $1/\sqrt{8}$ \\
      $\rho_s$ (spinorial) & 2 & 1 & 6 & $1/\sqrt{6}$ \\
      \hline
    \end{tabular}
  \end{center}

  The spinorial sector admits amplitude $1/\sqrt{6}$ versus $1/\sqrt{8}$ for the
  abelian sectors.
  The admissibility ratio is
  \begin{equation}
    R_p(Q_8) \;=\; \frac{A_{\rho_s}^{\max}}{A_{\rho_k}^{\max}}
    \;=\; \sqrt{\frac{8}{6}} \;=\; \sqrt{\frac{4}{3}} \;\approx\; 1.155.
  \end{equation}

  For $2I$ with its standard Cayley graph, an analogous calculation gives
  $\lambda_{\rho_s}=14$, $\lambda_{\rho_1}=9$ (lowest non-trivial abelian eigenvalue),
  yielding
  \begin{equation}
    R_p(2I) \;=\; \sqrt{\frac{14}{9}} \;\approx\; 1.247.
  \end{equation}

  The hierarchy is more pronounced in $2I$ than in $Q_8$, consistent with the
  strengthening of the spinorial selection described in Section~\ref{subsec:q8-to-2i}.
